% Standalone source for the VQC circuit figure (Model B).
% Edit this file, then recompile it to regenerate vqc_circuit.pdf:
%     pdflatex vqc_circuit.tex
% The main thesis includes the resulting PDF via \includegraphics.
\documentclass[border=8pt]{standalone}
\usepackage{tikz}
\usetikzlibrary{quantikz2}
\begin{document}
\begin{quantikz}[column sep=0.3cm, row sep=0.45cm]
\lstick{$\ket{0}$} & \gate{R_Y(x_1)}\gategroup[4,steps=6,style={dashed,rounded corners,inner xsep=4pt},background]{one layer, repeated $L=5$ times} & \gate{R_Y(\theta^{(\ell)}_1)} & \ctrl{1} &          &          & \targ{}   & \meter{} \\
\lstick{$\ket{0}$} & \gate{R_Y(x_2)} & \gate{R_Y(\theta^{(\ell)}_2)} & \targ{}   & \ctrl{1} &          &           & \meter{} \\
\lstick{$\ket{0}$} & \gate{R_Y(x_3)} & \gate{R_Y(\theta^{(\ell)}_3)} &           & \targ{}  & \ctrl{1} &           & \meter{} \\
\lstick{$\ket{0}$} & \gate{R_Y(x_4)} & \gate{R_Y(\theta^{(\ell)}_4)} &           &          & \targ{}  & \ctrl{-3} & \meter{}
\end{quantikz}
\end{document}
